% Options for packages loaded elsewhere
% Options for packages loaded elsewhere
\PassOptionsToPackage{unicode}{hyperref}
\PassOptionsToPackage{hyphens}{url}
\PassOptionsToPackage{dvipsnames,svgnames,x11names}{xcolor}
%
\documentclass[
  number]{elsarticle}
\usepackage{xcolor}
\usepackage{amsmath,amssymb}
\setcounter{secnumdepth}{5}
\usepackage{iftex}
\ifPDFTeX
  \usepackage[T1]{fontenc}
  \usepackage[utf8]{inputenc}
  \usepackage{textcomp} % provide euro and other symbols
\else % if luatex or xetex
  \usepackage{unicode-math} % this also loads fontspec
  \defaultfontfeatures{Scale=MatchLowercase}
  \defaultfontfeatures[\rmfamily]{Ligatures=TeX,Scale=1}
\fi
\usepackage{lmodern}
\ifPDFTeX\else
  % xetex/luatex font selection
\fi
% Use upquote if available, for straight quotes in verbatim environments
\IfFileExists{upquote.sty}{\usepackage{upquote}}{}
\IfFileExists{microtype.sty}{% use microtype if available
  \usepackage[]{microtype}
  \UseMicrotypeSet[protrusion]{basicmath} % disable protrusion for tt fonts
}{}
\makeatletter
\@ifundefined{KOMAClassName}{% if non-KOMA class
  \IfFileExists{parskip.sty}{%
    \usepackage{parskip}
  }{% else
    \setlength{\parindent}{0pt}
    \setlength{\parskip}{6pt plus 2pt minus 1pt}}
}{% if KOMA class
  \KOMAoptions{parskip=half}}
\makeatother
% Make \paragraph and \subparagraph free-standing
\makeatletter
\ifx\paragraph\undefined\else
  \let\oldparagraph\paragraph
  \renewcommand{\paragraph}{
    \@ifstar
      \xxxParagraphStar
      \xxxParagraphNoStar
  }
  \newcommand{\xxxParagraphStar}[1]{\oldparagraph*{#1}\mbox{}}
  \newcommand{\xxxParagraphNoStar}[1]{\oldparagraph{#1}\mbox{}}
\fi
\ifx\subparagraph\undefined\else
  \let\oldsubparagraph\subparagraph
  \renewcommand{\subparagraph}{
    \@ifstar
      \xxxSubParagraphStar
      \xxxSubParagraphNoStar
  }
  \newcommand{\xxxSubParagraphStar}[1]{\oldsubparagraph*{#1}\mbox{}}
  \newcommand{\xxxSubParagraphNoStar}[1]{\oldsubparagraph{#1}\mbox{}}
\fi
\makeatother


\usepackage{longtable,booktabs,array}
\newcounter{none} % for unnumbered tables
\usepackage{calc} % for calculating minipage widths
% Correct order of tables after \paragraph or \subparagraph
\usepackage{etoolbox}
\makeatletter
\patchcmd\longtable{\par}{\if@noskipsec\mbox{}\fi\par}{}{}
\makeatother
% Allow footnotes in longtable head/foot
\IfFileExists{footnotehyper.sty}{\usepackage{footnotehyper}}{\usepackage{footnote}}
\makesavenoteenv{longtable}
\usepackage{graphicx}
\makeatletter
\newsavebox\pandoc@box
\newcommand*\pandocbounded[1]{% scales image to fit in text height/width
  \sbox\pandoc@box{#1}%
  \Gscale@div\@tempa{\textheight}{\dimexpr\ht\pandoc@box+\dp\pandoc@box\relax}%
  \Gscale@div\@tempb{\linewidth}{\wd\pandoc@box}%
  \ifdim\@tempb\p@<\@tempa\p@\let\@tempa\@tempb\fi% select the smaller of both
  \ifdim\@tempa\p@<\p@\scalebox{\@tempa}{\usebox\pandoc@box}%
  \else\usebox{\pandoc@box}%
  \fi%
}
% Set default figure placement to htbp
\def\fps@figure{htbp}
\makeatother





\setlength{\emergencystretch}{3em} % prevent overfull lines

\providecommand{\tightlist}{%
  \setlength{\itemsep}{0pt}\setlength{\parskip}{0pt}}



 
\usepackage[]{natbib}
\bibliographystyle{elsarticle-num}


\makeatletter
\@ifpackageloaded{caption}{}{\usepackage{caption}}
\AtBeginDocument{%
\ifdefined\contentsname
  \renewcommand*\contentsname{Table of contents}
\else
  \newcommand\contentsname{Table of contents}
\fi
\ifdefined\listfigurename
  \renewcommand*\listfigurename{List of Figures}
\else
  \newcommand\listfigurename{List of Figures}
\fi
\ifdefined\listtablename
  \renewcommand*\listtablename{List of Tables}
\else
  \newcommand\listtablename{List of Tables}
\fi
\ifdefined\figurename
  \renewcommand*\figurename{Figure}
\else
  \newcommand\figurename{Figure}
\fi
\ifdefined\tablename
  \renewcommand*\tablename{Table}
\else
  \newcommand\tablename{Table}
\fi
}
\@ifpackageloaded{float}{}{\usepackage{float}}
\floatstyle{ruled}
\@ifundefined{c@chapter}{\newfloat{codelisting}{h}{lop}}{\newfloat{codelisting}{h}{lop}[chapter]}
\floatname{codelisting}{Listing}
\newcommand*\listoflistings{\listof{codelisting}{List of Listings}}
\makeatother
\makeatletter
\makeatother
\makeatletter
\@ifpackageloaded{caption}{}{\usepackage{caption}}
\@ifpackageloaded{subcaption}{}{\usepackage{subcaption}}
\makeatother
\usepackage{bookmark}
\IfFileExists{xurl.sty}{\usepackage{xurl}}{} % add URL line breaks if available
\urlstyle{same}
% fallback for those not using the hyperref driver hyperxmp:
\makeatletter
\@ifundefined{xmpquote}{\newcommand{\xmpquote}[1]{#1}}{}
\makeatother
\hypersetup{
  pdftitle={Assessing the influence of variable annual sampling effort on official national forest inventory statistics},
  pdfauthor={George C. Gaines III; Bryce Frank; KaDonna Randolph; Jim Westfall},
  pdfkeywords={\xmpquote{Forest Inventory and Analysis, national forest
inventory, sample design}},
  colorlinks=true,
  linkcolor={blue},
  filecolor={Maroon},
  citecolor={Blue},
  urlcolor={Blue},
  pdfcreator={LaTeX via pandoc}}


\setlength{\parindent}{6pt}
\begin{document}

\begin{frontmatter}
\title{Assessing the influence of variable annual sampling effort on
official national forest inventory statistics}
\author[1]{George C. Gaines III%
\corref{cor1}%
}
 \ead{george.gaines@usda.gov} 
\author[2]{Bryce Frank%
%
}

\author[3]{KaDonna Randolph%
%
}

\author[4]{Jim Westfall%
%
}


\affiliation[1]{organization={USDA Forest Service, Rocky Mountain
Research Station},city={Missoula},postcodesep={}}
\affiliation[2]{organization={USDA Forest Service, Pacific Northwest
Research Station},city={XXX},postcodesep={}}
\affiliation[3]{organization={USDA Forest Service, Southern Research
Station,},city={Knoxville},postcodesep={}}
\affiliation[4]{organization={USDA Forest Service, Northern Research
Station},city={XXX},postcodesep={}}

\cortext[cor1]{Corresponding author}




        
\begin{abstract}
In this study we\ldots{}
\end{abstract}





\begin{keyword}
    
    Forest Inventory and Analysis, national forest inventory, sample
design
\end{keyword}
\end{frontmatter}
    

\section{Introduction}\label{introduction}

The USFS Forest Inventory and Analysis (FIA) Program surveys and
analyzes the forests of the United States through a nationwide network
of permanent sample locations. These plot locations are distributed
across all lands with a spatial sampling intensity of about 1 location
per 2400 ha according to an unaligned, quasi-systematic equal
probability sample design. This design is spatially dispersed with
deviation from a regular grid to avoid alignment of sample locations
with periodic population features. Statistical inference on forest
population status and change parameters is enabled by measurements of
forest vegetation traits made at these locations by field crews.

Historically, FIA sample locations in each state were visited
periodically as time, budgets, and workforce capacity allowed. This
approach offered logistical flexibility but complicated analysis and
reporting timelines. In the early 2000's the program shifted from
periodic to annual inventory protocols calling for approximately equal
numbers of plot measurements each year on schedules intended to be fixed
and consistent within states. Under this annual inventory, each plot
location in a state is assigned to one of a set of panels, and all plot
locations assigned to a panel are designated for measurement in the same
year. The panels in a state together comprise an inventory cycle, and
inventory cycles range in length from 5 or 7 years in the eastern US to
10 years in the west. Once the final panel in an inventory cycle is
complete, remeasurements at sample locations in the first panel begin
the following year. For more information regarding the FIA sample design
see \citep[ Chapter 2]{gb1, gb2}.

In general, FIA requires that a panel be complete before the associated
tabular data are pushed to its public-facing databases, which the
program endeavors to do within 1 year of panel completion. Similarly,
FIA requires a complete inventory cycle within a state to publish
official state-scale estimates of forest parameters in fulfillment of
Congressional reporting obligations. In practice, however, adherence to
the annual design described above has proven very challenging or
impossible for the FIA Program. Instead, delayed and skipped field
measurements cause unintentional lengthening of inventory cycles and/or
survey nonresponse \citep[see, e.g.,][]{sarndal_nonresponse}, at rates
that vary regionally and over time.

The underlying causes of delayed and skipped field measurements are
varied, and can be chronic (e.g., denied access to private lands
\citep{westfall2022spatial}, logistical or administrative challenges in
remote regions \citep{affleck2026inventory}, persistent workforce and
budgetary constraints \citep{conroy_thesis}) or can result from discrete
events (e.g., the COVID19 pandemic \citep{covid}). Late measurements of
temporally dynamic variables can induce a ``time lag'' bias in point
estimates when the time component of the estimand is not adjusted
accordkingly. Further, survey nonresponse can introduce nonresponse bias
to estimates of population parameters when response propensity is
correlated with the attribute of interest. Better understanding of
patterns in late and skipped measurements \citep{conroy_thesis}, and the
investigation of new sampling and estimation techniques to ameliorate
the statistical consequences of nonresponse
\citep{westfall2022estimation, westfall2026double} are active research
topics, but beyond the scope of the present study.

Of immediate practical concern for FIA and its external partners are the
strain on field staff and delays to tabular data publication that result
from delayed field measurements. Methods have been proposed to control
the measurement delay process by disaggregating FIA's annual panel
structure into subpanels while still preserving the spatial dispersion
of the FIA sample network within field season years. Disaggregation of
panels into subpanels enables intentional alignment of workloads with
anticipated data acquisition workforce capacity in each field season
year through the application of variable annual field sampling effort
within individual states.

Though the logistical advantages of variable annual effort (VAE) are
profound (i.e., it can help ensure the long-term viability of the
inventory in the context of funding and staffing volatility), the
statistical consequences of the redistribution of partial panels within
inventory cycles, and of extended and variable inventory cycle lengths,
are unknown across states and epochs. In particular, a need exists to
understand and communicate to stakeholders how adoption of VAE protocols
may influence the statistical properties of forest parameter estimates
published by FIA in fulfillment of Congressional reporting imperatives.

Therefore, our research objectives in this study are to:

\begin{enumerate}
\def\labelenumi{\arabic{enumi}.}
\tightlist
\item
  define three VAE regimens with respect to subpanels (fast and slow
  start, dropped subpanels, and extended and shortened cycle lengths)
  intended to be representative of circumstances FIA could conceivably
  encounter in program administration
\item
  simulate the influence of those three VAE regimens on state-scale,
  inventory-cycle wide estimates of total net merchantable timber volume
  for three US states with distinct forest types, ecological attributes,
  and FIA inventory cycle lengths.
\item
  visualize and discuss simulation results in the context of
  implications for FIA program operations and the management of US
  forests.
\end{enumerate}

\section{Data \& Methods}\label{sec-data-methods}

\section{Discussion}\label{discussion}

\section{Conclusion}\label{conclusion}

\section*{References}\label{references}
\addcontentsline{toc}{section}{References}

\renewcommand{\bibsection}{}
\bibliography{vae.bib}





\end{document}
